%% 40장 특수함수 II
\chapter{Special Functions II: Incomplete Gamma and Beta Functions}
\label{ch:special-2}

\abstract{감마 분포, 카이제곱 분포, 포아송 분포의 누적분포함수는 모두 불완전 감마 함수이고, 베타, $t$, $F$, 이항 분포의 누적분포함수는 불완전 베타 함수이다. 통계 라이브러리의 정확도는 이 두 함수의 정확도에 달려 있다. 이 장에서는 불완전 감마 함수를 거듭제곱 급수와 연분수로 계산하는 방법을 유도하고, 두 방법이 모두 느려지고 소거가 생기는 전이 영역 $x\approx a$를 Temme의 균등 점근 전개가 메우는 방식을 본다. 두 방법에 공통으로 곱해지는 앞 인자 $x^ae^{-x}/\Gamma(a)$를 그대로 계산하면 큰 $a$에서 지수의 반올림이 결과를 망치는데, 이것을 Stirling 정규화로 푸는 방법을 정확한 계산과 견주어 보인다. 불완전 베타 함수는 DiDonato와 Morris의 다섯 가지 방법으로 영역을 나누어 계산하고 두 함수의 역함수는 도함수가 앞 인자를 $x$로 나눈 것이라는 사실을 이용한 Halley 반복으로 구한다.}

\section{The Incomplete Gamma Function}
\label{sec:incomplete-gamma}

\subsection{Definition}

\term{정규화된 불완전 감마 함수}{regularized incomplete gamma function}는
\[
P(a,x)=\frac1{\Gamma(a)}\int_0^xt^{a-1}e^{-t}\,dt,\qquad Q(a,x)=1-P(a,x)=\frac1{\Gamma(a)}\int_x^\infty t^{a-1}e^{-t}\,dt
\]
이다. $P(a,x)$는 모양 $a$의 감마 분포의 누적분포함수이고, 자유도 $\nu$의 카이제곱 분포는 $P(\nu/2,x/2)$, 평균 $\mu$인 포아송 분포의 $P(N\le k)$는 $Q(k+1,\mu)$이다. 39장의 오차함수도 특수한 경우로, $\operatorname{erf}x=P(1/2,x^2)$이다.

39장의 오차함수에서처럼 $P$와 $Q$는 각자 따로 계산해야 한다. 둘 가운데 작은 쪽을 정확하게 계산하고 큰 쪽은 $1$에서 빼는 것이 원칙이다. 작은 쪽을 $1$에서 빼면 작은 쪽의 낮은 자리 비트가 버려지지만, 결과가 $1/2$ 이상이므로 그 절대오차가 곧 상대오차로 작다. 대략 $x<a$이면 $P$가 작고, $x>a$이면 $Q$가 작다.

\subsection{Power Series}

\begin{proposition}
\label{prop:igamma-series}
$a>0$, $x\ge0$에서
\[
P(a,x)=\frac{x^ae^{-x}}{\Gamma(a+1)}\sum_{k=0}^\infty\frac{x^k}{(a+1)(a+2)\cdots(a+k)}
\]
이다.
\end{proposition}

\begin{proof}
$\gamma(a,x)=\int_0^xt^{a-1}e^{-t}dt$로 두고 부분적분을 하면 $\gamma(a,x)=x^ae^{-x}/a+\gamma(a+1,x)/a$이다. 이것을 $n$번 되풀이하면
\[
\gamma(a,x)=x^ae^{-x}\sum_{k=0}^{n-1}\frac{x^k}{a(a+1)\cdots(a+k)}+\frac{\gamma(a+n,x)}{a(a+1)\cdots(a+n-1)}
\]
이다. $\gamma(a+n,x)\le\int_0^xt^{a+n-1}dt=x^{a+n}/(a+n)$이므로 나머지 항은 $n\to\infty$에서 $0$으로 간다. $\Gamma(a)$로 나누고 $a\Gamma(a)=\Gamma(a+1)$을 쓰면 된다.
\end{proof}

급수의 항은 모두 양수이므로 소거가 없고 항의 비가 $x/(a+k)$이므로 $k>x-a$부터 기하급수적으로 준다. 그래서 $x<a$에서 $P$를 계산하는 데 알맞다.

\subsection{Continued Fractions}

$x>a$에서는 Legendre의 연분수
\[
\Gamma(a,x)=\frac{x^ae^{-x}}{x+1-a-\cfrac{1\cdot(1-a)}{x+3-a-\cfrac{2\cdot(2-a)}{x+5-a-\cdots}}}
\]
로 $Q$를 계산한다\cite{dlmf,cuyt2008cf}. 6장에서 본 대로 연분수는 앞에서부터 Lentz 방법으로 계산하거나, 필요한 항 수를 먼저 정한 뒤 꼬리부터 거꾸로 계산할 수 있다. 앞에서부터의 계산은 항 수를 미리 몰라도 되지만 걸음마다의 반올림이 곱해져 쌓인다. Gautschi는 불완전 감마 함수의 영역마다 급수와 연분수를 고르는 기준을 체계적으로 세웠고\cite{gautschi1979igamma}, DiDonato와 Morris의 알고리즘은 그 기준을 정밀하게 다듬었다\cite{didonato1986}.

\subsection{The Transition Region}

$x\approx a$이고 $a$가 크면 급수와 연분수가 모두 느리다. $x=a$에서 항(연분수는 보정 인자와 $1$의 차)이 부분합의 $10^{-17}$배 아래로 내려갈 때까지 binary64로 세어 보면, 급수는 $a=10,100,10^4,10^6$에서 $38$, $97$, $840$, $8012$항이 필요하고 연분수는 $10$, $41$, $198$, $987$항이 필요하다. 급수의 항 수는 대략 $\sqrt a$에, 연분수의 항 수는 그보다 느리게(대략 $a^{1/3}$) 늘어난다. 게다가 $P\approx Q\approx1/2$인 곳이어서 어느 쪽을 계산해도 상대 정확도를 잃을 여지가 있다. 감마 분포는 $a$가 크면 정규분포에 가까워지므로, 이 영역을 정규분포의 꼬리 확률로 근사하고 보정하는 것이 자연스럽다.

Temme은 이것을 $a$에 대해 균등한 점근 전개로 정확하게 만들었다\cite{temme1979igamma}. $\lambda=x/a$와
\[
\tfrac12\eta^2=\lambda-1-\log\lambda,\qquad\operatorname{sign}\eta=\operatorname{sign}(\lambda-1)
\]
로 $\eta$를 정의하면
\[
Q(a,x)=\tfrac12\operatorname{erfc}\bigl(\eta\sqrt{a/2}\bigr)+R_a(\eta),\qquad R_a(\eta)\sim\frac{e^{-a\eta^2/2}}{\sqrt{2\pi a}}\sum_{k\ge0}\frac{c_k(\eta)}{a^k}
\]
이다\cite{dlmf}. 계수 $c_k(\eta)$는 $\eta$의 해석 함수이고 $\eta=0$ 근처의 Taylor 계수로 계산한다. 첫 항이 정규분포의 꼬리를 정확한 변수 $\eta$로 쓴 것이고 나머지 항이 $a^{-1}$의 거듭제곱으로 줄어드는 보정이다. 이 전개는 $\lambda$가 $1$ 근처에 고정되지 않아도, 곧 꼬리 쪽에서도 $a$에 대해 균등하게 유효하다는 점에서 보통의 점근 전개보다 강하다. 계산에서 $\lambda-1-\log\lambda$는 $\lambda\approx1$에서 소거가 일어나므로 $\log(1+t)$의 급수나 \code{log1p}로 $t=\lambda-1$에서 직접 계산해야 한다.

\section{The Prefactor}
\label{sec:prefix}

\cref{prop:igamma-series}과 연분수와 Temme의 전개에는 모두 \term{앞 인자}{prefix} $x^ae^{-x}/\Gamma(a)$가 곱해진다. 이 값은 $P$의 도함수에 $x$를 곱한 것, 곧 $xP'(a,x)$이기도 하다. 정의대로 $\exp(a\log x-x-\log\Gamma(a))$로 계산하면, 지수 안의 세 항이 각각 $a\log a$ 정도의 크기여서 그 반올림 오차가 절대오차로 지수에 들어가고 결과의 상대오차가 된다. $a=1000$이면 세 항이 수천이므로 상대오차가 $10^{-13}$ 정도가 되고, $a$가 클수록 나빠진다.

Stirling 정규화가 이 문제를 푼다. $\Gamma^\ast(a)=\Gamma(a)/\bigl(\sqrt{2\pi/a}\,(a/e)^a\bigr)$로 두면 $\Gamma^\ast(a)\to1$이고, $\lambda=x/a$에 대해 $x^ae^{-x}=(a/e)^ae^{-a(\lambda-1-\log\lambda)}$이므로
\begin{equation}
\label{eq:prefix-stirling}
\frac{x^ae^{-x}}{\Gamma(a)}=\sqrt{\frac a{2\pi}}\;\frac{e^{-a\mu(\lambda-1)}}{\Gamma^\ast(a)},\qquad\mu(t)=t-\log(1+t)
\end{equation}
이다. 큰 부분 $(a/e)^a$가 정확히 상쇄되어 사라지고 지수에는 $a\mu(\lambda-1)$ 하나만 남는다. 이 값은 앞 인자의 로그와 부호가 반대이고 작은 항만큼만 다르므로, 결과가 아래로 넘치지 않는 한 크지 않다. $\mu(t)$는 $t$가 작으면 $t^2/2$처럼 작아서 \code{log1p}의 소거 없는 계산이 필요하다.

정확한 \code{decimal} 계산과 견주면 차이가 분명하다. $a=1000$, $x=1200$에서 정의대로의 계산은 상대오차가 $4.1\times10^{-13}$이고 \eqref{eq:prefix-stirling}을 binary64로 계산하면 $1.2\times10^{-14}$이다. $a=10^5$, $x=1.05\times10^5$에서는 $1.1\times10^{-10}$ 대 $1.2\times10^{-13}$이다. 뒤의 오차도 $a$와 함께 자란다. $\lambda=x/a$의 반올림이 $a\mu$에서 $|x-a|$배, 곧 $a(\lambda-1)$배로 커지기 때문이다. 그래서 정밀한 구현은 $a\mu(\lambda-1)$을 두 binary64의 합으로 계산한다.

\section{The Incomplete Beta Function}
\label{sec:incomplete-beta}

\term{정규화된 불완전 베타 함수}{regularized incomplete beta function}는
\[
I_x(a,b)=\frac1{B(a,b)}\int_0^xt^{a-1}(1-t)^{b-1}\,dt
\]
이고 대칭 $I_x(a,b)=1-I_{1-x}(b,a)$를 가진다. 이항 분포의 누적분포함수, $t$ 분포와 $F$ 분포의 누적분포함수가 모두 이 함수로 표현된다. 앞 인자는 $x^a(1-x)^b/(aB(a,b))$이고 불완전 감마와 같은 이유로 큰 $a$, $b$에서 정규화가 필요하다.

인자가 둘이라 영역을 나누는 일이 훨씬 복잡하다. 단순한 방법은 $x<(a+1)/(a+b+2)$이면 연분수로 $I_x(a,b)$를, 아니면 대칭으로 $I_{1-x}(b,a)$를 계산하는 것이다\cite{press2007nr}. 그러나 이 연분수는 $a$, $b$가 클 때 평균 근처에서 느리고 한 인자가 아주 작거나 클 때 정확도를 잃는다. DiDonato와 Morris의 TOMS 708은 영역마다 다섯 가지 방법을 고른다\cite{didonato1992}. 작은 $b$나 작은 $x$의 거듭제곱 급수, 인자를 정수만큼 옮기는 유한합, 큰 $a$와 작은 $b$에서 $x\approx1$의 점근 전개, 평균 아래의 연분수, 그리고 $a$, $b$가 모두 클 때 평균 근처의 Temme형 균등 점근 전개이다. 마지막 것은 불완전 감마 함수의 Temme 전개와 같은 생각으로, 이항 분포의 정규 근사를 정확하게 만든 것이다.

\section{Inverse Functions}
\label{sec:inverse-incomplete}

통계의 분위수 함수는 이 함수들의 역함수이다. 주어진 $p$에서 $P(a,x)=p$인 $x$를 찾는 일은 34장의 근 찾기인데, 도함수가 공짜라는 특별한 사정이 있다. $\partial P/\partial x=x^{a-1}e^{-x}/\Gamma(a)$는 바로 앞 인자를 $x$로 나눈 것이고 2차 도함수도 $\partial^2P/\partial x^2=\bigl((a-1)/x-1\bigr)\partial P/\partial x$로 바로 나온다. 그래서 3차 수렴하는 34장의 Halley 방법을 거의 추가 비용 없이 쓸 수 있다. 꼬리에서는 $P$가 몇 자릿수씩 변하므로 $P$ 대신 $\log P$에 대해 Newton 걸음을 하는 편이 낫다.

좋은 초기값도 중요하다. 카이제곱 분포에는 Wilson과 Hilferty의 근사, 곧 $(X/\nu)^{1/3}$이 평균 $1-2/(9\nu)$, 분산 $2/(9\nu)$인 정규분포에 가깝다는 사실을 쓴다\cite{wilson1931}. DiDonato와 Morris는 영역마다 초기값을 정교하게 골랐고\cite{didonato1986}, Gil, Segura, Temme은 균등 점근 전개를 거꾸로 써서 초기값을 만드는 방법을 정리했다\cite{gil2012igamma}. 불완전 베타 함수의 역함수에는 Cran, Martin, Thomas의 AS 109가 쓰인다\cite{cran1977}.

\section{Incomplete Gamma and Beta Functions in stardust}
\label{sec:special2-stardust}

\begin{instardust}[\code{docs/design/special.md} §5.3, §5.4]
stardust의 \code{gammainc}과 \code{gammaincc}는 $P$와 $Q$ 가운데 작은 쪽을 계산한다. $a=1$과 $a=1/2$은 닫힌 꼴이고, 작은 $x$와 작은 $a$는 Boost의 방법, 그 밖은 급수, 연분수, 그리고 $a>20$이고 $|x-a|<0.4a$이면 Temme 전개이다. Temme 전개의 계수 $c_k(\eta)$는 맞춘 값이 아니라 DLMF의 정의식을 정확한 유리수로 전개해 만들었다. 앞 인자는 \eqref{eq:prefix-stirling}의 지수 $a\mu(\lambda-1)$을 정확한 입력에서 double-double로 만들어 계산한다. 설계 문서는 scipy가 쓰는 Cephes가 $a\log x-x-\log\Gamma(a)$를 binary64로 계산해 꼬리에서 수천 ulp를 잃는다고 적는다. 연분수는 앞으로 가며 필요한 항 수를 찾고 값은 꼬리부터 거꾸로 계산하는데, 앞에서부터의 Lentz 방법은 60항에서 $17$ulp를 잃었다고 한다.
\end{instardust}

\begin{instardust}[\code{docs/design/special.md} §5.5]
\code{betainc}은 TOMS 708의 구조를 따르되 측정에 따라 가지를 조정했다. 두 인자가 모두 $1$보다 크고 작은 쪽이 $40$보다 작은 영역에서, TOMS 708은 유한합으로 $b$를 $1$ 이하로 내린 뒤 급수를 쓰지만 stardust는 연분수를 쓴다. 앞 인자를 세 번 대신 한 번 만들고 정확도도 약 두 배 좋았기 때문이다. 이항 분포의 누적분포함수의 최대 오차가 $10.1$ulp에서 $4.9$ulp로 줄었다고 설계 문서는 적는다. 역함수는 Halley 반복인데, 함숫값과 목표가 두 배 넘게 차이 나는 동안은 $\log F$에 대한 Newton 걸음을 쓴다.
\end{instardust}

\section{Concluding Remarks}

불완전 감마와 베타 함수의 계산은 39장에서 본 원칙들이 한꺼번에 시험받는 자리이다. 작은 쪽을 계산하고 소거 없는 표현을 고르며 큰 수의 지수를 정확히 상쇄하고 방법이 모두 느려지는 전이 영역은 균등 점근 전개로 메운다. 영역의 경계를 어디에 두느냐는 이론만으로 정해지지 않고 측정으로 다듬어진다. 다음 장에서는 진동하는 함수와 여러 매개변수를 가진 함수, 곧 Bessel 함수와 Airy 함수, 제타 함수, 초기하 함수로 넘어간다.

불완전 감마와 베타 함수의 계산 방법은 Gil, Segura, Temme의 책과 논문들\cite{gil2012igamma}과 DLMF의 8장\cite{dlmf}에 정리되어 있다.

\chapterreferences
